\subsection{All tame-local conjugacy classes in the actual inherited module quotient}

\begin{statement}{E.45}{Proved}{Complete local section coordinates, including zero first inertia}
Use the actual arithmetic group $H$ and local field
$K_{\mathfrak q}$ of E.42. Put
\[
 \mathscr A=\mathbb Z_3[[U]],\quad T=1+U,\quad
 \sigma=v_3(q-1),\quad
 \mathcal B_\sigma=\{0,3,6,\ldots,3^{\sigma+1}-3\}.
 \tag{E45.1}
\]
Every continuous local section, modulo geometric
conjugacy by $H$, has uniquely normalized four
coordinates
\[
 a,c\in3\mathbb Z_3,\qquad b,d\in\mathcal B_\sigma,
 \qquad
 s(\tau)=(a,v),\quad s(F)=(b,w),\quad v(1)=c,\quad w(1)=d.
 \tag{E45.2}
\]
After this normalization the following formulas
give all the classes, with no additional
identifications.

For $a\ne0$, the classes are exactly the
solutions of the finite two-row Smith system
(E42.6), with the representatives (E42.7)
and $h=h_0$. Every nonempty such set is the
explicit $\mathbb Z_3^{e-1}$ torsor of E.42,
where $e$ is the number of the affine orbits
(E42.4).

For $a=0$, put $r=b/(q-1)$. When
$r\in\mathbb Z_3$, every $c\in3\mathbb Z_3$
is allowed. The classes have the following
unique representatives, indexed by arbitrary
$z_{j,u}\in\mathbb Z_3$:
\[
 \begin{split}
 v&=cT^{-r},\\
 w&=d+T^{-r}
        \sum_{\substack{j\ge0\\
                     u\in(\mathbb Z/3^\sigma)^\times}}
                 z_{j,u}(T^{3^j u}-1).
 \end{split}
 \tag{E45.3}
\]
Choose the representatives $1\le u<3^\sigma$.
The series converges in the actual profinite
completed-group-ring topology, equivalently
the $(3,U)$-adic topology. Its parameter
space is the full compact product
\[
          \prod_{j\ge0}\prod_{u\in(\mathbb Z/3^\sigma)^\times}
                         \mathbb Z_3.
 \tag{E45.4}
\]

When $r\notin\mathbb Z_3$, put
$t=v_3(b)<\sigma$. A class exists exactly
for $c=0$. Its unique representatives,
indexed by $z_0,\ldots,z_{3^t-1}\in\mathbb Z_3$,
are
\[
 v=0,\qquad
 w=d+z_0q(1-T^b)
          +\sum_{i=1}^{3^t-1}z_i(T^i-1).
 \tag{E45.5}
\]
Thus this fiber is explicitly $\mathbb Z_3^{3^t}$.
Here $1\le t<\sigma$, since $b\in3\mathbb Z_3$.
The fixed-coordinate fibers in (E45.3) and
(E45.5) have their product and three-adic
topologies, respectively.

The precise linear maps underlying these
coordinates are also explicit. Set
\[
 D_b(f)=qT^bf(T^q)-f(T),\qquad f\in U\mathscr A.
 \tag{E45.6}
\]
For integral $r$, translation by $T^r$
gives
\[
 \operatorname{coker}(D_b:U\mathscr A\to U\mathscr A)
     \xrightarrow{\sim}
             \prod_{j,u}\mathbb Z_3.
 \tag{E45.7}
\]
For nonintegral $r$, the corresponding
map identifies this cokernel with the
actual lattice
\[
 \left\{(I_0,\ldots,I_{3^t-1})\in\mathbb Z_3^{3^t}:
                 \sum I_i\in(q-1)\mathbb Z_3\right\},
 \tag{E45.8}
\]
whose basis is $(q-1)e_0$ and
$e_i-e_0$, $1\le i<3^t$. The proof
defines the maps by weighted integrals
of integral measures and proves their
kernels and surjectivity.

\emph{Manuscript: \S51.16, folios733v--736v,
scans172--178, Part E(c),(d). This completes
all local section classes of the actual
quotient $H$ in E.38 and E.42. The maps
from the intrinsic metabelian and marked
groups, with their geometric kernels,
remain exactly (E38.5)--(E38.9).}
\end{statement}
\begin{proof}
\emph{The full conjugation action and normalization.}
All local cocycles are the actual generator
images satisfying the tame presentation
(E42.8). Conjugate their section by
$(s,u)\in H$, with
$s\in3\mathbb Z_3$ and $u(1)\in3\mathbb Z_3$.
Direct multiplication in the given
semidirect product gives
\[
\begin{aligned}
 a'&=a,&
 v'&=T^sv+(1-T^a)u,\\
 b'&=b+(1-q)s,&
 w'&=T^sw+u-qT^{b'}u(T^q),\\
 c'&=c,&d'&=d+(1-q)u(1).
\end{aligned}
 \tag{E45.9}
\]
In particular $a,c$ are actual invariants.
The two remaining augmentations vary
independently by $3(q-1)\mathbb Z_3$
$=3^{\sigma+1}\mathbb Z_3$.
Their representatives in
$\mathcal B_\sigma$ are unique. They
are achieved simultaneously by choosing
the unique required $s$ and taking $u$
to be the required constant with that
augmentation. Once $b,d$ have their
normalized values, (E45.9) preserves
them exactly when $s=0$ and $u(1)=0$.
The remaining conjugations are therefore
precisely $u\in U\mathscr A$. Formula
(E42.12) then proves the asserted
classification for $a\ne0$.

For $a=0$, E.42 gives $v=cT^{-r}$
when $r$ is integral, and gives
$v=0,c=0$ when $r$ is nonintegral.
The remaining $w$ is arbitrary of
augmentation $d$, and the residual
action is $w\mapsto w-D_b(u)$.
The constant $w=d$ is an actual
solution. Thus it remains to compute
the exact compact quotient
$U\mathscr A/D_b(U\mathscr A)$
and supply representatives. We do
that over the integral completed
group ring itself.

\emph{Integral measures and the free orbit calculation.}
Identify $\mathscr A$ with
$\mathbb Z_3[[\mathbb Z_3]]$ by
$T^x\leftrightarrow\delta_x$.
An element is an integral measure:
its coefficients on the finite
quotients $\mathbb Z/3^k$ give
compatible $\mathbb Z_3$ values
on clopen subsets. A continuous
$\mathbb Z_3$-valued function is
integrated by uniform limits of
its locally constant approximations.
The resulting integral exists
because all coefficients have
absolute value at most one.
Augmentation is total mass.
Under this identification,
$T^b\psi_q$, with
$\psi_q(f)=f(T^q)$, is pushforward
by the actual affine map
\[
                 \varphi(x)=qx+b.
 \tag{E45.10}
\]

On a compact free orbit
$\varphi^n(x_0)$, $n\in\mathbb Z_3$,
the parameter $n$ identifies its
measure module with
$\mathbb Z_3[[S-1]]$.
The operator $D_b$ becomes
multiplication by $qS-1$.
Evaluation at $S=q^{-1}$ gives
an exact continuous sequence
\[
 0\longrightarrow\mathbb Z_3[[S-1]]
   \xrightarrow{qS-1}\mathbb Z_3[[S-1]]
   \xrightarrow{I}\mathbb Z_3\longrightarrow0,
 \qquad
 I(\mu)=\int q^{-n}\,d\mu(n).
 \tag{E45.11}
\]
Indeed the power series may be
evaluated at $S-1=q^{-1}-1\in3\mathbb Z_3$;
division by the monic linear
factor proves that the kernel
is precisely $(qS-1)$. The
first map is injective in the
power-series domain. Dirac
mass at $n=0$ proves
surjectivity. This computes
the orbit cokernel as an
actual map of compact modules.

We shall use the elementary
valuation identity
\[
 v_3(q^n-1)=\sigma+v_3(n)
                 \quad(0\ne n\in\mathbb Z_3).
 \tag{E45.12}
\]
For positive integers $n$ it
follows by binomial expansion:
a unit exponent preserves the
first valuation, and cubing an
element of $1+3\mathbb Z_3$
raises it by exactly one.
The identity for all $n$ follows
by continuity on each annulus.
In particular $n\mapsto q^n$
is a homeomorphism from
$\mathbb Z_3$ onto
$1+3^\sigma\mathbb Z_3$.
Surjectivity can be checked
without a logarithm: (E45.12)
makes the normalized map a
bijection at every finite
three-power residue level,
and compactness passes to
the inverse limit.

\emph{Integral fixed point and infinitely many orbit coordinates.}
Suppose $r\in\mathbb Z_3$.
Multiplication by $T^r$ is an
automorphism of $U\mathscr A$,
and a direct substitution gives
\[
                T^rD_b(f)=D_0(T^rf).
 \tag{E45.13}
\]
We therefore compute $D_0$.
The punctured compact space
$\mathbb Z_3-\{0\}$ is the
disjoint union of clopen sets
\[
 A_{j,u}=3^j(u+3^\sigma\mathbb Z_3),
 \quad j\ge0,\quad
 u\in(\mathbb Z/3^\sigma)^\times.
 \tag{E45.14}
\]
By (E45.12) each is a single
free compact orbit, parametrized
by $x=3^juq^n$.
Its map (E45.11) is concretely
\[
 I_{j,u}(\mu)=\int_{A_{j,u}}\frac{3^ju}{x}\,d\mu(x).
 \tag{E45.15}
\]
The integrand takes integral
unit values and is continuous
on that clopen set.

There is an exact topological
module isomorphism
\[
 U\mathscr A\xrightarrow{\sim}
       \prod_{j,u}\mathscr A(A_{j,u}),\qquad
 \mu\longmapsto(\mu|_{A_{j,u}})_{j,u}.
 \tag{E45.16}
\]
Its inverse is the centered
measure series
\[
 (\mu_{j,u})\longmapsto
       \sum_{j,u}\bigl(\mu_{j,u}
                  -\operatorname{aug}(\mu_{j,u})\delta_0\bigr).
 \tag{E45.17}
\]
To check convergence, project
to the group ring of
$\mathbb Z/3^k$. Every term
with $j\ge k$ becomes zero
there, since its support and
its compensating mass both
reduce to zero. Only finitely
many terms remain. This
proves convergence and
continuity in the actual
profinite topology. It is
the $(3,U)$-adic topology;
no convergence in the
$U$-adic topology alone is
needed. The resulting measure
has total mass zero and the
prescribed restrictions.
For uniqueness, a measure
whose restrictions all vanish
is supported at zero: testing
any locally constant function
shows it equals its total
mass times $\delta_0$.
Total mass zero then makes
it zero.

The isomorphism (E45.16)
intertwines $D_0$ with its
orbit operators. Indeed
\[
 D_0(\nu-m\delta_0)
       =D_0\nu-(q-1)m\delta_0,
       \qquad\operatorname{aug}(D_0\nu)=(q-1)m.
\]
The product of the exact
sequences (E45.11) therefore
proves (E45.7), with no
closure or extra kernel.
The representative of the
coordinate $z_{j,u}$ is
$z_{j,u}(\delta_{3^ju}-\delta_0)$.
Transporting it back by
$T^{-r}$ and adding $d$
gives exactly (E45.3).
Each orbit integral of this
representative is its stated
coordinate. This proves both
surjectivity and uniqueness
of the section representatives.

\emph{No integral fixed point and the finite-rank lattice.}
Suppose $r\notin\mathbb Z_3$,
and put $t=v_3(b)<\sigma$.
The compact affine orbits are
now exactly
\[
 A_i=i+3^t\mathbb Z_3,
                  \qquad0\le i<3^t.
 \tag{E45.18}
\]
For a proof, the continuous
iterate is
\[
 \varphi^n(i)=q^ni+
                   b\frac{q^n-1}{q-1},\qquad
 \varphi^n(i)+r=q^n(i+r).
\]
Here $v_3(i+r)=t-\sigma<0$.
Equation (E45.12) yields
\[
 v_3(\varphi^n(i)-\varphi^m(i))
                         =t+v_3(n-m).
\]
Thus the normalized parameter
is an isometry on every
finite residue quotient and
a homeomorphism onto $A_i$.
Each orbit is free.

The finite partition into
these orbits and (E45.11)
give the explicit map
\[
 \mathscr A/D_b\mathscr A\xrightarrow{\sim}
     \mathbb Z_3^{3^t},\qquad
 I_i(\mu)=\int_{A_i}\frac{i+r}{x+r}\,d\mu(x).
 \tag{E45.19}
\]
The ratio is $q^{-n}$ under
the orbit parameter. It is
congruent to one modulo
$q-1$, so
\[
             \sum_i I_i(\mu)
                         \equiv\operatorname{aug}(\mu)
                                      \pmod{q-1}.
 \tag{E45.20}
\]
This sends $U\mathscr A$ into
the lattice (E45.8).

We prove the exact kernel
and image on $U\mathscr A$.
For any lattice vector
$I=(I_i)$, the measure
\[
 \mu_I=\sum_i I_i\delta_i
       -D_b\left(\frac{\sum_iI_i}{q-1}\delta_0\right)
 \tag{E45.21}
\]
has augmentation zero and
has weighted integrals $I_i$.
It proves surjectivity. If
all integrals of a measure
$\mu\in U\mathscr A$ vanish,
(E45.19) gives $\mu=D_b\nu$
in $\mathscr A$. Taking
augmentations yields
$0=(q-1)\operatorname{aug}(\nu)$,
so $\nu\in U\mathscr A$.
This is exactly the required
kernel, proving (E45.8).

The lattice basis
$(q-1)e_0,e_i-e_0$ has
representatives
\[
 q(1-T^b),\qquad T^i-1,
                     \quad1\le i<3^t.
\]
For the first, $\varphi(0)=b$
and its weighted integral is
$q(1-q^{-1})=q-1$ on orbit
zero; for the others the
integrals are $e_i-e_0$.
This proves (E45.5), including
its uniqueness. All quotient
maps above are continuous
surjections between compact
Hausdorff modules, and the
displayed representative
maps are continuous. Hence
the asserted topologies also
follow from these actual maps.

\emph{Finite certificates and their transition maps.}
The program
\texttt{e\_tame\_classes.py}
works on the exact finite
group-ring lattices
$\mathbb Z_3[\mathbb Z/3^k]$
in the augmentation basis
$\delta_j-\delta_0$.
For an affine cycle of
length $\ell$, the full
coinvariant is
$\mathbb Z_3/(q^\ell-1)$
by the weighted coordinate
$\sum_nq^{-n}a_n$.
If the cycle exponents
$e_i=\sigma+v_3(\ell_i)$
are ordered increasingly,
the augmentation-kernel
coinvariants have exponents
\[
 e_1-\sigma,e_2,\ldots,e_f,
 \tag{E45.22}
\]
with a zero exponent omitted.
This follows directly from
the kernel of the sum map
to $\mathbb Z_3/(q-1)$:
use the first coordinate
multiplied by $q-1$ and
the differences from that
coordinate. The program
independently verifies the
Smith matrices and this
formula in $48$ cases.

It also verifies $36$ actual
transition maps from level
$k+1$ to $k$, rather than
only comparing invariant
factors. If a higher cycle
starts at a point whose
projection has lower orbit
parameter $n_0$, its weighted
coordinate contributes
$q^{-n_0}$ to that lower
coordinate, modulo the lower
cycle's $q^\ell-1$.
The program checks this
formula on every Dirac
mass, and checks that
reduction of the group
ring commutes with $D_b$.
The proposed representatives
in (E45.3) and (E45.5)
are independently verified
to generate every finite
augmentation coinvariant.
Three complete modular
Smith transformation
certificates are retained
in the JSON output.
The measure proofs above
establish the all-level
isomorphisms themselves.

Finally the actual quotient
maps of E.38 send every
intrinsic local section to
one of the classes just
computed. Their specified
geometric kernels remain
part of this comparison.
The additional intrinsic
lifting tests are E.37
and E.44, with their
genuine arithmetic actions;
the computation here
completes exactly the
local section set of $H$.
\end{proof}
